% preamble-exam-zh.tex — 共享 preamble
% 用于所有 exam-zh 模板
%
% 样式策略：
%   屏幕 → 语义彩色（蓝=关键想法，红=易错，绿=路线，灰=提示/教师备注）
%   打印 → 自动降级为黑白（通过 print mode 或 monochrome 选项）

\documentclass{exam-zh}

% ---------- 基础配置 ----------
\examsetup{
  page/size = a4paper,
  page/show-head = false,
  page/show-foot = false,
  sealline/show = false,
  font = times,
  math-font = xits,
}

\usepackage{tcolorbox}
\usepackage{needspace}
\tcbuselibrary{skins,breakable}

% ---------- 打印/屏幕颜色切换 ----------
% 用法：\documentclass[monochrome]{exam-zh} 即可切换为纯黑白
% 注意：不要使用 \if@xxx 放在 makeatother 外部；这里改成普通条件名。
\newif\ifeduprintcolor
\eduprintcolortrue

\makeatletter
\@ifclasswith{exam-zh}{monochrome}{\eduprintcolorfalse}{}
\makeatother

% ---------- 语义颜色定义 ----------
% 屏幕：有色彩区分；打印：降级为灰度
\ifeduprintcolor
  % --- 屏幕彩色模式 ---
  \definecolor{edu-blue}{RGB}{47, 82, 122}       % 关键想法
  \definecolor{edu-blue-bg}{RGB}{243, 246, 252}
  \definecolor{edu-red}{RGB}{180, 40, 40}         % 易错提醒
  \definecolor{edu-red-bg}{RGB}{255, 245, 240}
  \definecolor{edu-green}{RGB}{34, 100, 60}       % 路线图
  \definecolor{edu-green-bg}{RGB}{242, 250, 245}
  \definecolor{edu-orange}{RGB}{160, 90, 0}       % 训练目标
  \definecolor{edu-orange-bg}{RGB}{255, 248, 235}
  \definecolor{edu-gray}{RGB}{100, 100, 100}      % 提示/教师备注
  \definecolor{edu-gray-bg}{RGB}{248, 248, 248}
  \definecolor{edu-purple}{RGB}{90, 50, 120}      % 思考题
  \definecolor{edu-purple-bg}{RGB}{248, 244, 252}
\else
  % --- 打印黑白模式 ---
  \definecolor{edu-blue}{gray}{0.15}
  \definecolor{edu-blue-bg}{gray}{0.96}
  \definecolor{edu-red}{gray}{0.25}
  \definecolor{edu-red-bg}{gray}{0.97}
  \definecolor{edu-green}{gray}{0.20}
  \definecolor{edu-green-bg}{gray}{0.97}
  \definecolor{edu-orange}{gray}{0.25}
  \definecolor{edu-orange-bg}{gray}{0.98}
  \definecolor{edu-gray}{gray}{0.40}
  \definecolor{edu-gray-bg}{gray}{0.97}
  \definecolor{edu-purple}{gray}{0.30}
  \definecolor{edu-purple-bg}{gray}{0.98}
\fi

% ---------- 自定义教学环境 ----------
% 共性：enhanced + breakable（跨页不断裂）

% 原题卡片 — 强边框，突出显示
\newtcolorbox{problemcard}{
  enhanced, breakable,
  colback=edu-blue-bg,
  colframe=edu-blue,
  boxrule=1.2pt,
  arc=0pt,
  left=3mm, right=3mm, top=3mm, bottom=3mm,
}

% 解题路线图 — 左边线 + 浅底
\newtcolorbox{routebox}{
  enhanced, breakable,
  colback=edu-green-bg,
  colframe=edu-green!30,
  boxrule=0pt,
  borderline west={2.5pt}{0pt}{edu-green},
  left=3mm, right=3mm, top=3mm, bottom=3mm,
}

% 关键想法 — 蓝色左边线
\newtcolorbox{keyideabox}{
  enhanced, breakable,
  colback=edu-blue-bg,
  colframe=edu-blue!30,
  boxrule=0pt,
  borderline west={2.5pt}{0pt}{edu-blue},
  left=3mm, right=3mm, top=3mm, bottom=3mm,
}

% 易错提醒 — 红色左边线
\newtcolorbox{mistakebox}{
  enhanced, breakable,
  colback=edu-red-bg,
  colframe=edu-red!20,
  boxrule=0pt,
  borderline west={3pt}{0pt}{edu-red},
  left=3mm, right=3mm, top=2.5mm, bottom=2.5mm,
}

% 提示 — 灰色虚线左边线
\newtcolorbox{hintbox}{
  enhanced, breakable,
  colback=edu-gray-bg,
  colframe=edu-gray!20,
  boxrule=0pt,
  borderline west={2pt}{0pt}{edu-gray, dashed},
  left=3mm, right=3mm, top=2mm, bottom=2mm,
}

% 思考题 — 紫色虚线左边线
\newtcolorbox{thinkbox}{
  enhanced, breakable,
  colback=edu-purple-bg,
  colframe=edu-purple!20,
  boxrule=0pt,
  borderline west={2pt}{0pt}{edu-purple, dashed},
  left=2.5mm, right=2.5mm, top=2.5mm, bottom=2.5mm,
}

% 教师备注 — 灰色左边线，小字
\newtcolorbox{teachernote}{
  enhanced, breakable,
  colback=edu-gray-bg,
  colframe=edu-gray!15,
  boxrule=0pt,
  borderline west={2pt}{0pt}{edu-gray!60},
  left=2.5mm, right=2.5mm, top=2.5mm, bottom=2.5mm,
  fontupper=\small,
  colupper=black!60,
}

% 训练目标 — 橙色左边线，小字
\newtcolorbox{traininggoal}{
  enhanced, breakable,
  colback=edu-orange-bg,
  colframe=edu-orange!15,
  boxrule=0pt,
  borderline west={1.5pt}{0pt}{edu-orange},
  left=2mm, right=2mm, top=1mm, bottom=1mm,
  fontupper=\small,
}

% 答题区 — 无边框，纯留白
\newcommand{\answerarea}[1][40mm]{%
  \vspace{2mm}%
  \begin{tcolorbox}[
    enhanced, breakable,
    colback=white,
    colframe=white,
    boxrule=0pt,
    height=#1,
    valign=top,
    bottom=0mm,
    after skip=0mm,
  ]
  \end{tcolorbox}%
}

% ---------- 字体设置 ----------
% exam-zh (ctexbook) already sets CJK fonts via ctex.
% Only override if specific fonts are needed.
% macOS: Songti SC, STHeiti, STFangsong
% Linux: SimSun, SimHei, FangSong

% ---------- 页面设置 ----------
% exam-zh (ctexbook) already loads geometry.
% Override margins if needed:
% \geometry{top=16mm, bottom=16mm, left=14mm, right=14mm}

\examsetup{
  question/show-answer = true,
  fillin/show-answer = true,
  solution/show-solution = show-stay,
}

% ---------- 教师版解答样式 ----------
\newtcolorbox{solutionblock}{
  enhanced, breakable,
  colback=edu-blue-bg,
  colframe=edu-blue!25,
  boxrule=0pt,
  borderline west={2pt}{0pt}{edu-blue!50},
  arc=0pt,
  left=3mm, right=3mm, top=2mm, bottom=2mm,
  before skip=2mm,
  after skip=3mm,
  fontupper=\small,
}

% 解答步骤编号
\newcounter{solstep}
\newcommand{\solstep}[2]{%
  \stepcounter{solstep}%
  \par\medskip
  \noindent\hangindent=6mm
  {\color{edu-blue}\bfseries\textcircled{\scriptsize\thesolstep}\enspace #1}\enspace #2\par
}

\begin{document}

\section*{三垂直与旋转等腰直角专题（四） · 练习卷（教师版）}

\section*{一、选择题（每题 12 分，共 48 分）}


\begin{keyideabox}
  \textbf{中等思想：共端点旋转（手拉手）}
  手拉手模型的本质是：两个等腰直角（或正方形/等边三角形）共顶点。通过 SAS 证明两翼三角形全等，并由此得出对应线段相等与旋转 $90^\circ$ （或 $60^\circ$）的垂直/夹角关系。
\end{keyideabox}



\begin{question}[points=12]
  以 $\triangle ABC$ 的边 $AC, AB$ 向外作正方形 $ACNM$ 和 $ABEF$，连接 $BM, CF$。则 $\triangle FAC$ 与 $\triangle BAM$ 全等的判定依据是
  \begin{choices}
    \item SAS
    \item ASA
    \item AAS
    \item SSS
  \end{choices}
  \paren[A]
  \begin{solutionblock}
    $AB = AF$, $AC = AM$。夹角 $\angle FAC = \angle BAM = 90^\circ + \angle BAC$。判定依据为边角边 (SAS)。
  \end{solutionblock}
\end{question}



\begin{question}[points=12]
  接上题，线段 $BM$ 与 $CF$ 的几何位置关系是
  \begin{choices}
    \item 互相垂直且相等
    \item 平行且相等
    \item 只相等不垂直
    \item 相交夹角为 $45^\circ$
  \end{choices}
  \paren[A]
  \begin{solutionblock}
    由全等 $\triangle FAC \cong \triangle BAM$ 得 $BM = CF$。且该全等相当于绕点 $A$ 旋转 $90^\circ$ 的几何变换，对应线段夹角亦为 $90^\circ$，故 $BM \perp CF$。
  \end{solutionblock}
\end{question}



\begin{question}[points=12]
  已知正方形向外作时，若 $A(0,0)$, $B(3,0)$，且 $C$ 在第一象限。正方形 $ABEF$ 向外作，则点 $F$ 的纵坐标为
  \begin{choices}
    \item -3
    \item 3
    \item 0
    \item -1.5
  \end{choices}
  \paren[A]
  \begin{solutionblock}
    因为正方形向外作，$C$ 在 $x$ 轴上方，则正方形 $ABEF$ 必须伸向 $x$ 轴下方，即第四象限。边长为 $3$，故 $F(0, -3)$。纵坐标为 $-3$。
  \end{solutionblock}
\end{question}



\begin{question}[points=12]
  已知 $F(0, -3)$， $M\left(-\dfrac{\sqrt{15}}{2}, -\dfrac{1}{2}\right)$。则线段 $FM$ 的长度为
  \begin{choices}
    \item $\sqrt{10}$
    \item 10
    \item 5
    \item $\sqrt{5}$
  \end{choices}
  \paren[A]
  \begin{solutionblock}
    $FM^2 = \left(-\frac{\sqrt{15}}{2} - 0\right)^2 + \left(-\\ frac{1}{2} - (-3)\right)^2 = \frac{15}{4} + \left(\frac{5}{2}\right)^2 =\ \ \frac{15}{4} + \frac{25}{4} = \frac{40}{4} = 10 \implies FM = \sqrt{10}$。
  \end{solutionblock}
\end{question}


\section*{二、填空题（每题 12 分，共 12 分）}


\begin{question}[points=12]
  以 $\triangle ABC$ 的边 $AB, AC$ 向外作等边 $\triangle ABD$ 和等边 $\triangle ACE$，连接 $CD, BE$。则线段 $CD$ 与 $BE$ 的数量关系为 \fillin。
  \paren[$CD = BE$]
  \begin{solutionblock}
    $AD = AB$, $AC = AE$。夹角 $\angle DAC = \angle BAE = 60^\circ + \angle BAC$。由 SAS 得 $\triangle DAC \cong \triangle BAE \implies CD = BE$。
  \end{solutionblock}
\end{question}


\section*{三、解答题（共 40 分）}

\needspace{8\baselineskip}

\begin{problem}[points=20]
  以 $\triangle ABC$ 的边 $AC$、点 $AB$ 向外作正方形 $ABEF$ 和正方形 $ACNM$，联结 $BM$、点 $CF$。

(1) 证明： $BM = CF$，且 $BM \perp CF$；
(2) 若 $AB = 3$，$AC = 2$，$BC = 4$，建立平面直角坐标系求出线段 $FM$ 的长度。
  \answerarea[65mm]
  \setcounter{solstep}{0}
  \begin{solutionblock}
    \textbf{答案：}(1) 证明 SAS 全等与旋转 90 度；(2) $FM = \sqrt{10}$\par\medskip
    \solstep{证明两线段相等}{$\because$ 四边形 $ABEF$ 与 $ACNM$ 是正方形。\\ $\therefore AB = AF$，$AC = AM$，且 $\angle FAB = \angle CAM = 90^\circ$。\\ $\therefore \angle FAC = 90^\circ + \angle BAC = \angle BAM$。\\ 在 $\triangle FAC$ 和 $\triangle BAM$ 中，有 $AF = AB$，$\angle FAC = \angle BAM$，$AC = AM$。\\ $\therefore \triangle FAC \cong \triangle BAM$ (SAS)。\\ $\therefore BM = CF$。}
    \solstep{证明两线段垂直}{因为 $\triangle FAC \cong \triangle BAM$。这相当于将 $\triangle FAC$ 绕点 $A$ 顺时针旋转 $90^\circ$ 得到 $\triangle BAM$。\\ 旋转变换中，对应线段 $CF$ 与 $BM$ 的夹角等于旋转角。\\ 故它们的夹角为 $90^\circ$，即 $BM \perp CF$。}
    \solstep{建系求 FM 长度}{设 $A(0,0)$， $AB$ 落在 $x$ 轴正半轴， 故 $B(3,0)$。\\ 设 $C(x_c, y_c)$。由 $AC=2,\ \ BC=4$ 得：\\ $\begin{cases} x_c^2 + y_c^2 = 4 \\ (x_c-3)^2 + y_c^2 =\ \ 16 \end{cases} \implies x_c^2 - 6x_c + 9 + y_c^2 = 16 \implies 4 - 6x_c\ \ + 9 = 16 \implies 6x_c = -3 \implies x_c = -1/2$。\\ 代入得 $y_c = \sqrt{4\ \ - 1/4} = \frac{\sqrt{15}}{2}$ (第一象限)。\\ 因为正方形向外作：$F(0, -3)$。\\ $AM\ \ = AC = 2$， 且 $AM \perp AC$ $\implies M$ 坐标通过 $C$ 顺时针旋转 $90^\circ$ 得到：\\ \ $M\left(-\frac{\sqrt{15}}{2}, -1/2\right)$。\\ $FM = \sqrt{\left(-\\ frac{\sqrt{15}}{2}\right)^2 + \left(-\frac{1}{2} - (-3)\right)^2} = \\ sqrt{\frac{15}{4} + \frac{25}{4}} = \sqrt{10}$。}
  \end{solutionblock}
\end{problem}


\needspace{8\baselineskip}

\begin{problem}[points=20]
  以 $\triangle ABC$ 的边 $AB, AC$ 向外作等边 $\triangle ABD$ 和等边 $\triangle ACE$，连接 $CD, BE$。

(1) 证明： $CD = BE$；
(2) 证明线段 $CD$ 与线段 $BE$ 的交角为 $60^\circ$。
  \answerarea[60mm]
  \setcounter{solstep}{0}
  \begin{solutionblock}
    \textbf{答案：}(1) SAS全等证明；(2) 夹角为 60 度，由旋转 60 度几何性质保证\par\medskip
    \solstep{证明全等线段相等}{$\because \triangle ABD$ 与 $\triangle ACE$ 是等边三角形。\\ $\therefore AD = AB$，$AC = AE$， 且 $\angle DAB = \angle CAE = 60^\circ$。\\ $\therefore \angle DAC = 60^\circ + \angle BAC = \angle BAE$。\\ 在 $\triangle DAC$ 和 $\triangle BAE$ 中，有 $AD = AB$，$\angle DAC = \angle BAE$，$AC = AE$。\\ $\dots \triangle DAC \cong \triangle BAE$ (SAS) $\implies CD = BE$。}
    \solstep{证明两线夹角}{因为 $\triangle DAC \cong \triangle BAE$。这相当于将 $\triangle DAC$ 绕点 $A$ 逆时针旋转 $60^\circ$ 得到 $\triangle BAE$。\\ 旋转变换中，对应线段 $CD$ 和 $BE$ 的夹角等于旋转角。\\ 故它们的交角为 $60^\circ$。}
  \end{solutionblock}
\end{problem}


\clearpage
\section*{参考答案}


\begin{keyideabox}
  \textbf{选择题}
  1. A  2. A  3. A  4. A
\end{keyideabox}



\begin{keyideabox}
  \textbf{填空题}
  1. CD = BE
\end{keyideabox}



\begin{keyideabox}
  \textbf{解答题}
  第 1 题：(1) 证明 SAS 全等与旋转 90 度；(2) $FM = \sqrt{10}$
第 2 题：(1) SAS全等证明；(2) 夹角为 60 度，由旋转 60 度几何性质保证
\end{keyideabox}



\end{document}
